The simulation campaign of \cref{sec:simulation} establishes a single architectural commitment: that for cooperative aerial manipulation in which the structural transition self-announces through actuator-side measurements on the plant's native timescale, the announcement-detection-reconfiguration timeline that defines classical active fault-tolerant pipelines is unnecessary, and a single decentralized feed-forward identity~$T_i^{\mathrm{ff}} = T_i$ in the per-drone altitude channel suffices. The five-drone, ten-kilogram dual-fault scenario is a single point on the actuator-margin feasibility surface (\cref{fig:feasibility-envelope}). In this section, we discuss what the campaign establishes, how it changes the design space, and where its boundary leaves open architectural questions that lie beyond the present work.

\subsection{Audit of the Three Contributions}
\label{sec:discussion-audit}

The taut-cable reduction theorem of \cref{thm:reduction} (C1) is the precondition on which the entire stability argument rests, and its empirical defense has two layers. The domain-gate audit of \cref{tab:domain-audit} certifies that all six pre-declared missions sit inside the admissibility set $\Omega^{\mathrm{dwell}}_{\tau}$ on every gate, with worst-case slack-run duration $\SI{35.8}{\milli\second}$, slack duty cycle $1.66\,\%$, and QP active-set transition fraction $0.10\,\%$ against thresholds $\SI{40}{\milli\second}$, $2.5\,\%$, and $1.0\,\%$ respectively --- non-trivial margin on every gate, on every mission. The reduction-fidelity layer of \cref{sec:VI-reduction} then quantifies, on the same trajectories, the per-rope tension dispersion that the lumped scalar abstracts: cruise-window RMS residual sits at $\SI{12.4}{\newton}$ on dual-fault variants, with peaks reaching $\SI{60}{\newton}$ at the most centripetally-loaded instants (\cref{tab:reduction-residuals}). The dispersion is informative rather than alarming. The reduction's $\mathcal{O}(\delta+\eta_{\max})$ guarantee bounds the closed-loop trajectory deviation between the full plant and the reduced plant under the same controller, not the per-rope tension difference, and the feed-forward identity reads each drone's \emph{own} measured tension --- so the dispersion of $T_i$ about the lumped mean has no actuator-channel pathway into the altitude error. What remains absent is a direct trajectory-level comparison between the bead-chain plant and a parallel reduced-rope simulator under the same controller; this is the closure that would meet the reduction theorem's own standard rather than the standard of its domain-gate hypotheses, and it is identified as the principal extension of the C1 audit.

The hybrid practical-ISS recovery envelope of \cref{thm:hybrid-stability} (C2) is supported through three independent measurements that converge on the same conclusion. The integrated system-level performance of \cref{tab:performance-ci} shows that all six variants pass the pre-registered acceptance thresholds with the dual-fault peak altitude sag at $\SI{95.2}{\milli\metre}$ against the $\SI{100}{\milli\metre}$ ceiling. Direct fits on the V4 and V5 Lyapunov proxy $V(t)=\xi^{\top}P_v\xi$ yield empirical contraction coefficients $\hat{\rho}_{\mathrm{V4}}=0.20$ and $\hat{\rho}_{\mathrm{V5}}=0.045$, both well below the $\rho<1$ sufficiency bound of Lemma~\ref{lem:dwell-decay}. The dwell-time boundary characterisation of \cref{sec:VI-dwell} extends this further: $\hat{\rho}<1$ holds across the swept range $\tau_d/\tau_{\mathrm{pend}}\in[0.5,3.0]$ including the sub-threshold point $\tau_d=0.5\,\tau_{\mathrm{pend}}$, where $\hat{\rho}\approx 0.63$. The H2 hypothesis is therefore sufficient but not necessary for empirical contraction on the demonstrated trajectory family --- a finding that documents margin within the analytical envelope without claiming any guarantee outside it. The leading-order analytical estimate $\rho_{\mathrm{ana}}\approx e^{-\alpha_{xy}\tau_{\mathrm{pend}}} \approx 0.005$ is two orders of magnitude tighter than the measured contraction; the gap reflects the post-fault equilibrium shift and the finite-time integration contributions that the leading-order treatment neglects, and the measured contraction is the appropriate quantity to anchor design margin against.

The closed-form $L_1$ sampling-rate bound of \cref{prop:gamma-star} (C3) is the most narrowly scoped of the three contributions and accordingly receives the most directly targeted evidence. The mass-mismatch sweep of \cref{sec:VI-E} shows the $L_1$ layer recovering approximately $40\,\%$ of the FF-disabled altitude sag across the $56\,\%$ mass variation of P2-B, confirming the layer's design-regime function. The adaptive-gain stability dashboard of \cref{fig:l1-stability-dashboard} then validates \cref{prop:gamma-star} directly: across $\Gamma\in [5\!\times\!10^{2},\,8\!\times\!10^{4}]$ spanning $2.4$ decades and reaching $1.7\,\Gamma^{\star}$, the empirical altitude RMSE is unimodal in $\Gamma$ with minimum at $\Gamma\approx 10^{4}$, the peak sag follows the same shape, and the late-vs-early variance ratio remains within $[0.36,\,0.37]$ across the full sweep. The smooth performance degradation past $\Gamma^{\star}$, rather than a sharp instability transition, is precisely what a sufficient-condition Lyapunov bound predicts: $\Gamma^{\star}$ identifies the regime where contraction is \emph{guaranteed} analytically; the empirical practical-stability boundary on this trajectory family lies beyond it.

The three contributions converge on a single causal anchor. The P2-A ablation campaign holds every other controller parameter fixed and disables the identity $T_i^{\mathrm{ff}}=T_i$ in the QP layer (\cref{tab:ff-ablation}). The cruise RMSE rises by $34$--$39\,\%$ and the peak post-fault payload altitude sag inflates by $3.6$--$4.0\times$ across V3, V4, and V5; the directional consistency across all three fault schedules and the four-row self-announcement causal chain of \cref{fig:self-announcement} --- which traces the response from measured tension through feed-forward command, commanded thrust, and altitude error within a single $2.5\,\si{\second}$ window --- isolate the identity as the mechanism on which C1, C2, and C3 jointly depend. Without it, the altitude channel sees the post-fault tension redistribution as an unrejected matched disturbance, the bounded fault-jump of Lemma~\ref{lem:jump-continuity} inflates by the corresponding overshoot, and the dwell-cycle decay of Lemma~\ref{lem:dwell-decay} cannot absorb the resulting transient within one pendulum period. The ablation makes this concrete in measurable units.

\subsection{What the Architecture Changes}
\label{sec:discussion-advance}

Existing decentralized cooperative-transport architectures address the fault-tolerance question through one of three families. Active fault-tolerant pipelines~\cite{LiuSuspensionFailure2021} estimate a tension-related disturbance, classify it as a fault when a detection threshold is crossed, and reassign load through a scheduler; the detection-confirmation-reassignment timeline of this loop is, in agile multirotor regimes, comparable to or longer than the payload-pendulum period $\tau_{\mathrm{pend}}\approx \SI{2.24}{\second}$ on which the load redistribution itself unfolds. Passive force-control and leader--follower architectures~\cite{GassnerPassiveForce2017,NoCommTransport2021} avoid the inter-UAV exchange but presuppose a fixed cable-cardinality regime, with no closed-loop guarantee across the constraint-set transition. Centralized and synchronous-distributed force allocation~\cite{Michael2011,WangOng2017,LiLoianno2023,PallarLi2025,SuBhowmick2023} is structurally inapplicable, because the per-tick global-state synchronization on which it depends is precisely what the constraint-cardinality transition breaks at the actuator-tick rate.

The architectural commitment of the present work is to read the structural transition off the actuator-side measurement on the plant's native timescale and to convert it directly into a thrust increment through the single decentralized identity $T_i^{\mathrm{ff}}=T_i$, eliminating the announcement-detection-reconfiguration timeline as a system-level construct rather than as a faster substitute for it. The empirical signature of that commitment is the per-fault recovery times recorded in \cref{tab:recovery-iae}: $\SI{0.00}{\second}$ (first faults of V3, V4, V5, V6, sub-tick) and $\SI{0.32}{\second}$ and $\SI{0.33}{\second}$ on the second faults of V4 and V6 respectively, both well inside one pendulum period and measured against the $\sim\!\SI{2}{\second}$ timeline that any detection-based architecture would consume on threshold confirmation alone. The thrust-overshoot dynamics of \cref{tab:thrust-transient} make the same statement on the actuator channel: surviving drones spike to $1.18$--$1.79\!\times$ their post-fault settled thrust within one second of each cable severance, decay back to a new (elevated) steady state within roughly one payload-pendulum period, and never approach the $0.9\,f_{\max}$ saturation hazard line of \cref{tab:actuator-margin}.

The four architectural properties stated in the introduction follow as consequences of the commitment, audited individually by the campaign. \emph{Locality} is enforced by construction: the per-drone admissible information set $\mathcal{I}_i$ defined in \cref{eq:info-set} contains no peer state, no peer-cable tension, no fault flag, and no shared scheduler. \emph{Passivity} --- the absence of detector dynamics, classification logic, and reassignment rule from the closed loop --- is the architectural face of the same identity: the same equation that operates pre-fault produces the post-fault response, and the self-announcement chain of \cref{fig:self-announcement} traces the entire causal path through measurement, feed-forward, and thrust without an intermediate logical step. \emph{Robustness to multi-agent failure} is established by induction on the fault counter: \cref{thm:hybrid-stability}'s recovery envelope carries unchanged at each event, the architectural cost of an additional failure is one more bounded Lyapunov jump $\chi(\Delta_f)$ that the identity controls, and the closed-loop spectrum exponentially contracts that jump within one inter-fault dwell. The campaign demonstrates the $F=1$ and $F=2$ instances of this induction; the theorem extends to $F\le N-2$ on the actuator-margin condition, and the demonstrated configuration consumes $27\,\%$ of the actuator envelope (\cref{fig:feasibility-envelope}), leaving designable margin in $(N,m_L,F)$. \emph{Team-size independence} is the structural counterpart: the per-drone law and the lemma stack of \cref{sec:stability} carry no $N$ structurally, and the dependence concentrates entirely in the two scalar quantities $m_Lg/N$ and $m_Lg/(N{-}F)$. The certificate is therefore invariant in $N$ within the admissibility envelope; the five-drone case is a single point on a designable feasibility surface, not a case-study identity.

The two null findings reported in the campaign are read as architectural evidence for this framing rather than as failed extension claims. The MPC tension-ceiling layer (P2-C) shows essentially identical peak rope tension $\approx \SI{104.6}{\newton}$ across all tested ceilings $T_{\max}\in[60,100]\,\si{\newton}$ (\cref{tab:mpc-ceiling-sweep}), because the trajectory-induced pre-fault spike at the climb-to-lemniscate transition dominates the post-fault quasi-static load that the ceiling targets. The reshape supervisor (P2-D) returns $0.0\,\%$ measurable benefit across $T_{\mathrm{ref}}\in\{6,8,10,12\}\,\si{\second}$ (\cref{tab:reshape-period-sweep}, \cref{fig:reshape-binding-T6}), because the dynamic centripetal load $a_c\in[0.82,3.29]\,\si{\metre\per\second\squared}$ modulates the per-rope tension distribution throughout the trajectory and dominates the hover-equilibrium asymmetry that the linearised closed-form assignment of \cref{thm:reshape} optimises against. Both nulls demonstrate that, on the demonstrated trajectory family, the conventional active-fault-tolerance machinery is \emph{structurally inactive}, not because it has been disabled but because the underlying feed-forward identity has already satisfied the relevant constraint at the actuator-margin and per-rope-load levels at which those layers would otherwise become binding. This is the empirical face of the claim that detection and reconfiguration are not the mechanisms doing the work, and the bookkeeping argument of \cref{sec:theorem-preservation} ensures that enabling the layers outside their binding regimes preserves the recovery envelope of \cref{thm:hybrid-stability} unchanged.

\subsection{Boundary of the Architectural Category}
\label{sec:discussion-boundary}

The campaign characterizes its operating envelope through three classes of evidence --- domain-gate satisfaction, internal margin, and adversarial probes --- and each class identifies a distinct boundary at which the present analysis stops. Within the analysis domain of \cref{thm:hybrid-stability}, three regimes lie outside the swept envelope. The rapid-multi-fault regime, with $\tau_d<\tau_{\mathrm{pend}}$, voids the analytical contraction guarantee of Lemma~\ref{lem:dwell-decay}; the empirical sweep of \cref{sec:VI-dwell} shows that the controller delivers contraction margin even at $\tau_d=0.5\,\tau_{\mathrm{pend}}$ on the canonical trajectory family, but a parameter regime that empirically violates contraction --- indicatively heavier payload combined with sub-threshold dwell to push the slack-excursion budget into $\eta>\eta_{\max}$ --- has not been identified, and its construction is the natural completion of the dwell-boundary characterisation. The long-slack regime would push the trajectory outside $\Omega^{\mathrm{dwell}}_{\tau}$ and disable the tension-cancellation mechanism of Lemma~\ref{lem:tension-cancel}; the QP-fallback regime, in which solver infeasibility forces the component-wise clamp of \cref{sec:qp-and-ff}, would break Doctrine D1. Both are open to the same operating-boundary search, and both would benefit from a multi-fault extension of the dwell-time argument under a controlled fault-schedule density that generalizes Lemma~\ref{lem:dwell-decay} beyond pairwise composition.

A second class of extensions targets the linearisation residuals already characterized in the present analysis. The tube-MPC tightening that folds the $\Gamma_f e^{-\lambda(t-t_k^{\star})}$ overshoot of the recovery envelope of \cref{thm:hybrid-stability} into the constraint right-hand side of the tick QP \eqref{eq:mpc-qp} is the natural strengthening of the C3-MPC layer for spike-dominated regimes, and would convert the P2-C null into a positive performance result on trajectory families for which the dynamic spike rather than the quasi-static load is the binding constraint. An aggressive-period reshape study with an objective extended to include the dynamic centripetal load --- rather than the quasi-static-hover target of \cref{thm:reshape} --- would similarly close the P2-D null. Main-mode $L_1$ evaluation at the canonical $\SI{10}{\kilo\gram}$ operating point with an updated gain-selection from \cref{prop:gamma-star} would complete the C3-$L_1$ robustness picture along the payload axis. A four-axis robustness response surface over $(m_L,k_s,c_{\mathrm{drag}},\bar{w})$ on $n=100$ CRN-paired seeds is the visualization counterpart of the same robustness claim and is independent of any controller modification.

A third class of conditions lies outside the architectural category itself; these --- drone-side faults, soft cable failures, tension-sensor failures, and payload-attitude dynamics --- were enumerated in \cref{sec:theorem-preservation} as out-of-category and require a different announcement mechanism. They are not gaps in the present campaign.

Within this architectural category, the simulation campaign of \cref{sec:simulation} establishes that a single decentralized feed-forward identity~$T_i^{\mathrm{ff}}=T_i$ is sufficient to deliver hybrid practical input-to-state stability with measured contraction margin across single and dual unannounced cable-severance events, that the identity is necessary in the precise sense the P2-A ablation isolates, and that the conventional active fault-tolerance machinery is structurally inactive on the trajectory family for which the identity is binding.
